%In recent years, Graph Neural Networks (GNNs) have emerged as a powerful framework for anomaly detection by identifying fraudulent nodes whose behaviors differ from those of benign nodes \citep{GADBench, GADSurvey}. Existing GNN-based anomaly detection approaches can be classified into two major categories: spectral methods and spatial methods. 
\noindent \textbf{Spatial GNNs for GAD.}
Spatial methods typically detect anomalies by aggregating and transforming information from local neighborhoods. 
To overcome the smoothing effects in spatial filters, most methods focus on improving neighbor selection and
aggregation under fraud camouflage and graph heterophily. 
For example, CARE-GNN \citep{CARE-GNN} employs reinforcement learning to adaptively select informative neighbors and suppress potentially irrelevant messages. 
PC-GNN \citep{PCGNN} further addresses fraud camouflage and severe class imbalance through label-balanced node sampling and selective neighborhood aggregation. 
Additionally, GAGA~\citep{GAGA} improves anomaly detection on low-homophily graphs by organizing neighboring nodes into label-aware groups and encoding their group-level interactions. 
In contrast,  ConsisGAD~\citep{ConsisGAD} focuses on limited supervision by combining pseudo-label generation with consistency regularization.
%to exploit the information contained in unlabeled nodes.

\noindent \textbf{Spectral GNNs for GAD.}
Spectral GNNs are an important branch of GNNs that interpret node features as graph signals and process them through frequency-domain operations~\citep{xu2020graph}.
Early approaches typically adopt basic linear filters, such as GCN~\citep{GCN}, which is a first-order approximation of ChebyNet with low-pass behavior.
Several later approaches employ polynomial parameterizations to construct more flexible spectral responses. 
AMNet ~\citep{AMNet} learns separate low- and high-frequency filters based on Bernstein polynomials. Similarly, Jacobi \citep{JacobiConv} and BernNet \citep{BernNet} construct adaptive filters using Jacobi and Bernstein polynomial bases to adapt the filter response to different spectral regions.  
Despite their different bases and parameterizations, these methods mainly focus on learning a suitable frequency response over a given graph operator.
Another line of research develops graph-wavelet-based models which jointly capture spatial locality and frequency information for GAD tasks.
BWGNN \citep{BWGNN} identifies a spectral right-shift phenomenon associated with anomalous nodes and constructs a Beta-wavelet filter bank to capture anomaly-related information across multiple frequency bands.
GHRN~\citep{GHRN} exploits the relationship between graph heterophily and high-frequency signals and suppresses potentially harmful inter-class connections to improve anomaly detection.
Building upon multi-frequency spectral modeling, recent detectors further diversify filter
designs, such as hybrid spectral filters~\citep{SEC-GFD},  trainable dynamic wavelets~\citep{DSGAD} and Chi-Square wavelet filters~\citep{ChiGAD}.

\noindent \textbf{Signed Operators and Edge Heterophily.}
Classical spectral analysis of signed graphs shows that edges carrying negative weights yield a signed Laplacian that is positive semi-definite with valid spectrum~\citep{KunegisSigned}.
MSGNN~\citep{MSGNN} later extends this construction to a magnetic signed Laplacian for directed signed graphs.
In the heterophilic setting, FAGCN~\citep{FAGCN} learns signed attention coefficients to mix low- and high-frequency signals.
H\textsuperscript{2}-FDetector~\citep{H2-FDetector} propagates opposite-sign messages along predicted heterophilic connections.
Recently, GREET~\citep{GREET} proposes a method to explicitly discriminate
homophilic from heterophilic edges to guide representation learning instead of pruning them.


